\documentclass[11pt]{article}

\usepackage[letterpaper,margin=0.9in]{geometry}
\usepackage[T1]{fontenc}
\usepackage[utf8]{inputenc}
\usepackage{lmodern}
\usepackage{amsmath,amssymb,amsfonts,mathtools}
\usepackage{graphicx}
\usepackage{booktabs}
\usepackage{array}
\usepackage{enumitem}
\usepackage{fancyhdr}
\usepackage{xcolor}
\usepackage{hyperref}
\usepackage{float}
\usepackage{titlesec}
\usepackage{tikz}
\usepackage{longtable}
\usepackage{multicol}
\usepackage{caption}
\usepackage{setspace}
\usepackage{bm}

\definecolor{CCDdark}{RGB}{17,57,96}
\definecolor{CCDaccent}{RGB}{217,119,6}
\definecolor{CCDgreen}{RGB}{47,133,90}
\definecolor{CCDlight}{RGB}{234,242,248}
\definecolor{CCDsoft}{RGB}{244,244,245}
\definecolor{CCDgray}{RGB}{90,100,112}
\definecolor{CCDpurple}{RGB}{124,58,237}
\definecolor{CCDred}{RGB}{178,58,72}

\hypersetup{
  colorlinks=true,
  linkcolor=CCDdark,
  urlcolor=CCDdark,
  citecolor=CCDdark,
  pdftitle={Chapter 6: Optimal Control and Closed-Loop CCD},
  pdfauthor={Saeid Bayat}
}

\pagestyle{fancy}
\fancyhf{}
\fancyhead[L]{\small\color{CCDgray}Control Co-Design}
\fancyhead[R]{\small\color{CCDgray}Chapter 6: Optimal Control and Closed-Loop CCD}
\fancyfoot[C]{\thepage}
\setlength{\headheight}{14pt}

\titleformat{\section}{\Large\bfseries\color{CCDdark}}{\thesection}{0.5em}{}
\titleformat{\subsection}{\large\bfseries\color{CCDdark}}{\thesubsection}{0.5em}{}
\titleformat{\subsubsection}{\normalsize\bfseries\color{CCDdark}}{\thesubsubsection}{0.5em}{}

\newcommand{\vect}[1]{\boldsymbol{#1}}
\newcommand{\xp}{\vect{x}_p}
\newcommand{\xc}{\vect{x}_c}
\newcommand{\xstate}{\vect{x}}
\newcommand{\uvec}{\vect{u}}
\newcommand{\yvec}{\vect{y}}
\newcommand{\wvec}{\vect{w}}
\newcommand{\R}{\mathbb{R}}

\newenvironment{learningobjectives}{%
\begin{quote}\textbf{Learning objectives.} After completing this chapter, you should be able to:
\begin{enumerate}[leftmargin=1.4em,itemsep=0.35em,topsep=0.5em]
}{\end{enumerate}\end{quote}}

\newcommand{\keyideabox}[1]{%
\begin{center}
\fcolorbox{CCDgreen}{white}{\begin{minipage}{0.92\linewidth}\small #1\end{minipage}}
\end{center}}

\begin{document}

% Cover
\thispagestyle{empty}
\begin{tikzpicture}[remember picture,overlay]
  \fill[CCDdark] (current page.north west) rectangle ([yshift=-2.6in]current page.north east);
  \fill[CCDaccent] ([yshift=-2.65in]current page.north west) rectangle ([yshift=-2.75in]current page.north east);
  \fill[blue!8] (current page.south west) rectangle ([yshift=0.35in]current page.south east);
\end{tikzpicture}

\vspace*{0.7in}
{\color{white}\fontsize{18}{20}\selectfont CONTROL CO-DESIGN\par}
\vspace{0.18in}
{\color{white}\Large Integrated Optimization of Physical and Control Systems\par}

\vspace{2.15in}
{\fontsize{28}{32}\selectfont\bfseries\color{CCDdark} Chapter 6\par}
\vspace{0.18in}
\noindent\begin{minipage}{0.92\linewidth}
\raggedright
{\fontsize{22}{26}\selectfont\bfseries\color{black!85} Optimal Control and Closed-Loop CCD\par}
\end{minipage}
\vspace{0.20in}
\par\noindent{\large\color{CCDgray} From ideal open-loop trajectories to implementable feedback controllers\par}

\vspace{0.55in}
\noindent\fcolorbox{CCDgreen}{white}{%
\begin{minipage}{0.84\linewidth}
\large The key distinction in this chapter is information. An open-loop optimizer assumes a planned future, while a closed-loop controller reacts to measurements as the future unfolds.
\end{minipage}}

\vfill
{\large Saeid Bayat\par}
\vspace{0.12in}
{\normalsize Draft course chapter for senior undergraduate and beginning graduate students\par}
\vspace{0.12in}
{\normalsize July 2026\par}
\newpage

\tableofcontents
\newpage

\section*{Chapter Overview}
\addcontentsline{toc}{section}{Chapter Overview}

In earlier chapters, we formulated general control co-design (CCD) problems and discussed several possible solution architectures. We now turn to an important modeling issue inside those architectures: \emph{What kind of control object is being optimized?}

Sometimes the optimizer chooses a full control history $\uvec(t)$ over a finite time interval. This is the classical setting of \textbf{optimal control}. Sometimes the optimizer chooses a finite set of controller parameters, such as gains in a feedback law. These two viewpoints are related, but they are not the same. The first often gives an ideal performance benchmark; the second usually leads to controllers that are actually implementable in practice.

This chapter develops that distinction carefully. We begin with open-loop optimal control, then compare control trajectories with controller parameters, then move to feedback-controller co-design and model predictive control (MPC). A recurring theme is \textbf{information availability}: what the controller knows now, what it can predict, and what it cannot know exactly in advance.

\begin{learningobjectives}
  \item explain the idea of open-loop optimal control and its role in CCD;
  \item distinguish between optimizing a control trajectory and optimizing controller parameters;
  \item explain why feedback-controller co-design is important for implementable engineering systems;
  \item discuss the limitations of assuming perfect future information;
  \item identify different kinds of information available to a controller during operation;
  \item explain the basic idea of model predictive control;
  \item describe how open-loop optimal-control solutions can be turned into implementable closed-loop controllers.
\end{learningobjectives}

\section{Why This Chapter Matters}

When students first see an optimal-control problem, it can appear to solve everything at once. After all, if we can optimize the control signal over time, why worry about controller structure at all?

The answer is that a control trajectory and a controller are different objects:
\begin{itemize}[leftmargin=1.4em]
  \item a \textbf{control trajectory} is a time function planned in advance;
  \item a \textbf{controller} is a rule that maps available information into actions during operation.
\end{itemize}

In a real engineering system, disturbances, noise, and uncertainty make it impossible to know the exact future. That is why practical systems almost always rely on feedback. Figure~\ref{fig:ch6-map} summarizes the main ideas of the chapter.

\begin{figure}[H]
\centering
\includegraphics[width=0.98\textwidth]{figures/fig_chp6_1.png}
\caption{Concept map for the chapter. The main thread runs from open-loop control to implementable closed-loop control.}
\label{fig:ch6-map}
\end{figure}

\keyideabox{\textbf{Main idea.} Open-loop optimal control often gives a useful ideal benchmark or design insight, but most engineered systems must ultimately use a feedback controller that responds to measured information during operation.}

\section{Open-Loop Optimal Control}

\subsection{Basic definition}

Open-loop optimal control (OLOC) optimizes a control function over a time interval, typically written as $\uvec(t)$ for $t\in[t_0,t_f]$. The control is computed using the model, the initial condition, and any assumed future information, and is then treated as a planned signal.

A generic continuous-time optimal-control problem can be written as
\begin{align}
\min_{\uvec(\cdot),\xstate(\cdot)} \quad & J = \Phi\big(\xstate(t_f)\big) + \int_{t_0}^{t_f} L\big(\xstate(t),\uvec(t),t\big)\,dt \\
\text{subject to} \quad & \dot{\xstate}(t) = f\big(\xstate(t),\uvec(t),t\big), \\
& \xstate(t_0)=\xstate_0, \\
& \vect{c}\big(\xstate(t),\uvec(t),t\big)\le 0, \\
& \vect{b}\big(\xstate(t_f)\big)\le 0.
\end{align}

Here, the optimizer chooses the time functions $\uvec(t)$ and $\xstate(t)$ so as to minimize the objective while satisfying the system dynamics and constraints.

\subsection{Interpretation}

The word \emph{open-loop} means the control does not change in response to newly measured state information during the planned horizon. The control sequence is computed first; then it is applied as planned.

\begin{figure}[H]
\centering
\includegraphics[width=0.92\textwidth]{figures/fig_chp6_2.png}
\caption{Open-loop optimal control. The control trajectory is planned using the model, the initial condition, and assumed future information.}
\label{fig:ch6-oloc}
\end{figure}

This is not necessarily wrong. In some applications, an open-loop plan is entirely reasonable. For example:
\begin{itemize}[leftmargin=1.4em]
  \item if the horizon is very short,
  \item if uncertainty is small,
  \item if the trajectory is only a benchmark or design study,
  \item if the optimized control is recomputed repeatedly, as in MPC.
\end{itemize}

\subsection{Role in control co-design}

In CCD, OLOC is often used to answer questions such as:
\begin{itemize}[leftmargin=1.4em]
  \item If the plant were designed differently, what is the best achievable open-loop performance?
  \item How much energy can ideally be extracted or saved?
  \item What qualitative features appear in the best control action?
  \item What plant characteristics make the ideal control task easier or harder?
\end{itemize}

Thus, OLOC is valuable even if it is not the final implementation method.

\subsection{Example}

Consider a wave-energy device. One can optimize the PTO force trajectory over a fixed wave record. The result may reveal an upper bound on energy capture under that assumed wave history. That does not automatically give an implementable controller, but it provides a useful benchmark and design insight.

\section{Control Trajectories versus Controller Parameters}

\subsection{Two different kinds of decision variables}

In many optimization problems, students use the word ``control'' to describe whatever control-related quantity is being optimized. But two very different choices are possible:
\begin{enumerate}[leftmargin=1.5em]
  \item optimize the time-varying signal itself, $\uvec(t)$ or its discretized samples;
  \item optimize parameters of a controller law, such as gains in $\uvec=\kappa(\yvec;\xc)$.
\end{enumerate}

Figure~\ref{fig:ch6-traj-vs-param} compares these two representations.

\begin{figure}[H]
\centering
\includegraphics[width=0.95\textwidth]{figures/fig_chp6_3.png}
\caption{Trajectory optimization versus controller-parameter optimization. The first optimizes a signal directly; the second optimizes a policy or law that generates the signal from information.}
\label{fig:ch6-traj-vs-param}
\end{figure}

\subsection{Optimizing control trajectories}

When a control trajectory is optimized, the decision variables may be the control values at mesh points:
\begin{equation}
\uvec_1,\uvec_2,\ldots,\uvec_N.
\end{equation}

This gives the optimizer great flexibility. It can shape the control signal almost arbitrarily, subject to bounds and the dynamic model.

Advantages include:
\begin{itemize}[leftmargin=1.4em]
  \item high flexibility,
  \item strong ability to reveal best-case performance,
  \item a natural formulation for optimal-control theory.
\end{itemize}

Limitations include:
\begin{itemize}[leftmargin=1.4em]
  \item the result may be specific to one scenario or disturbance realization,
  \item the control may be hard to implement directly,
  \item no explicit feedback law is produced.
\end{itemize}

\subsection{Optimizing controller parameters}

Alternatively, one may choose a controller structure and optimize its parameters:
\begin{equation}
\uvec(t)=\kappa\big(\yvec(t);\xc\big),
\end{equation}
where $\yvec(t)$ denotes the available measurements and $\xc$ is a vector of gains or coefficients.

Examples include:
\begin{itemize}[leftmargin=1.4em]
  \item PID gains,
  \item state-feedback gains,
  \item observer or estimator parameters,
  \item basis-function weights,
  \item tuning weights in an MPC cost function.
\end{itemize}

Advantages include:
\begin{itemize}[leftmargin=1.4em]
  \item direct implementability,
  \item explicit feedback structure,
  \item usually lower-dimensional decision variables.
\end{itemize}

Limitations include:
\begin{itemize}[leftmargin=1.4em]
  \item the controller structure may be restrictive,
  \item some ideal performance may be lost relative to full trajectory optimization,
  \item choosing a suitable controller structure is itself a design decision.
\end{itemize}

\subsection{Why the distinction matters in CCD}

A plant designed under ideal open-loop control may not remain optimal when only a realistic feedback controller is allowed. Therefore, the plant design can depend strongly on which control representation is used.

\section{Feedback-Controller Co-Design}

\subsection{Basic idea}

Feedback-controller co-design optimizes plant variables together with parameters of a feedback controller. The controller uses measured information during operation and therefore can react to disturbances or uncertainty.

\begin{figure}[H]
\centering
\includegraphics[width=0.92\textwidth]{figures/fig_chp6_6.png}
\caption{Feedback-controller co-design. The controller is represented by parameters in a realizable feedback law, and performance is assessed in closed loop.}
\label{fig:ch6-feedback-ccd}
\end{figure}

A generic formulation is:
\begin{align}
\min_{\xp,\xc} \quad & J\big(\xp,\xc\big) \\
\text{subject to} \quad & \dot{\xstate}(t)=f\big(\xstate(t),\uvec(t),\xp,\wvec(t)\big), \\
& \uvec(t)=\kappa\big(\yvec(t);\xc\big), \\
& \yvec(t)=h\big(\xstate(t),\xp\big), \\
& \text{state, control, and design constraints.}
\end{align}

Here, the control signal is generated by a feedback law rather than optimized independently at every time instant.

\subsection{What changes compared with OLOC}

Compared with open-loop optimal control, the decision variables are no longer the entire future control history. Instead, the optimizer chooses the parameters $\xc$ of a feedback policy. The dynamic response is then determined by the interaction of:
\begin{itemize}[leftmargin=1.4em]
  \item the plant design $\xp$,
  \item the controller parameters $\xc$,
  \item the realized disturbances,
  \item the available measurements.
\end{itemize}

\subsection{Why this is more realistic}

A feedback controller can reject disturbances, reduce sensitivity to model mismatch, and respond to state deviations as they occur. For this reason, most real engineering systems---vehicles, robots, wind turbines, marine devices, and process systems---operate in closed loop.

\subsection{Example}

In wind-turbine control co-design, one may optimize tower or structural parameters together with pitch-controller gains. The measured rotor speed and other signals are used during operation, allowing the controller to respond to changing wind conditions.

\section{Limitations of Perfect Future Information}

\subsection{The hidden assumption in many open-loop problems}

An OLOC formulation often assumes that the future disturbance history or reference signal is known exactly over the horizon. For example, in a simulation study, one may optimize a control trajectory for a known wave record or a known road profile.

That is mathematically convenient, but it is often physically unrealistic.

\begin{figure}[H]
\centering
\includegraphics[width=0.95\textwidth]{figures/fig_chp6_4.png}
\caption{Perfect future information is rarely available. Some information is known now, some is forecast with uncertainty, and some is fundamentally unknown in advance.}
\label{fig:ch6-perfect-future}
\end{figure}

\subsection{Why this matters}

If the future is assumed to be known exactly, the optimizer may exploit details that a real controller could never know during operation. The resulting performance can therefore be too optimistic.

This does not make OLOC useless. It simply means we must interpret the result correctly:
\begin{itemize}[leftmargin=1.4em]
  \item as an ideal upper bound,
  \item as a benchmark for comparison,
  \item as a source of design insight,
  \item not automatically as a realizable control law.
\end{itemize}

\subsection{Illustrative example}

Suppose an active suspension controller is optimized using a road profile that is assumed known perfectly over the next 10 seconds. If the actual road deviates from that assumed profile, the open-loop sequence may perform poorly. A feedback controller, by contrast, can respond to the actual motion and load transfer observed in real time.

\subsection{Distinguishing forecast from certainty}

In some applications, useful forecasts do exist. For instance:
\begin{itemize}[leftmargin=1.4em]
  \item weather forecasts for wind or wave energy,
  \item route previews for autonomous vehicles,
  \item upcoming demand estimates in process systems.
\end{itemize}

However, forecasts are not the same as exact future knowledge. A controller must usually account for forecast error, not ignore it.

\section{Information Availability}

\subsection{A control problem is an information problem}

A good way to think about controller design is to ask: \emph{What information is available at the moment the control action must be chosen?}

This question is central to both control and CCD because the plant may be redesigned to make useful information more available or more informative.

\begin{figure}[H]
\centering
\includegraphics[width=0.88\textwidth]{figures/fig_chp6_5.png}
\caption{Information availability. A controller can only use information that is measured, estimated, or predicted with realistic uncertainty.}
\label{fig:ch6-info}
\end{figure}

\subsection{Three categories of information}

It is useful to separate information into three categories:
\begin{enumerate}[leftmargin=1.5em]
  \item \textbf{Available now:} measured states, outputs, or commands available at the current time;
  \item \textbf{Estimated or predicted:} forecasts, observers, estimators, or preview signals;
  \item \textbf{Unavailable exactly:} future disturbance realizations, future sensor noise, and modeling errors.
\end{enumerate}

\subsection{CCD implications}

In control co-design, plant design can affect information availability. For example:
\begin{itemize}[leftmargin=1.4em]
  \item adding sensors or changing sensor placement improves observability;
  \item mechanical redesign may make important dynamics easier to measure;
  \item actuator changes may influence what control authority is achievable;
  \item structural design may reduce sensitivity to uncertain disturbances.
\end{itemize}

Thus, closed-loop CCD often involves not just plant and controller parameters, but also the information pattern available to the controller.

\keyideabox{\textbf{Important viewpoint.} A control architecture should not assume more information than a real system can obtain. A large part of practical controller design is choosing actions based on the right information set, not based on impossible foresight.}

\section{Model Predictive Control}

\subsection{Core idea}

Model predictive control (MPC) is one of the most important bridges between open-loop optimal control and practical feedback control. At each time step, MPC solves a finite-horizon optimal-control problem using the current state estimate and possibly forecasts over a short future horizon. It then applies only the first control move and repeats the process at the next time step.

\begin{figure}[H]
\centering
\includegraphics[width=0.95\textwidth]{figures/fig_chp6_7.png}
\caption{Receding-horizon idea in model predictive control. Each new optimization produces a local open-loop plan, but only the first action is implemented before the plan is updated.}
\label{fig:ch6-mpc}
\end{figure}

\subsection{Why MPC is a feedback method}

Although MPC repeatedly solves an open-loop optimization problem, it is still a feedback controller because each new optimization uses the most recent measurements or state estimates. The optimization horizon moves forward in time, which is why MPC is also called \emph{receding-horizon control}.

\subsection{Generic structure}

At time step $k$:
\begin{enumerate}[leftmargin=1.5em]
  \item measure or estimate the current state $\xstate_k$;
  \item solve a finite-horizon optimal-control problem;
  \item obtain an optimized sequence $\uvec_k,\uvec_{k+1},\ldots,\uvec_{k+N-1}$;
  \item apply only the first control action $\uvec_k$;
  \item advance to the next step and repeat.
\end{enumerate}

\subsection{Advantages of MPC}

MPC is attractive because it:
\begin{itemize}[leftmargin=1.4em]
  \item incorporates a model directly,
  \item handles multivariable systems naturally,
  \item accommodates state and control constraints,
  \item uses updated measurements, making it closed loop,
  \item can incorporate short-term forecasts or preview information.
\end{itemize}

\subsection{Limitations of MPC}

MPC also has challenges:
\begin{itemize}[leftmargin=1.4em]
  \item an optimization problem must be solved online,
  \item the controller depends on model quality,
  \item the chosen horizon length and discretization matter,
  \item real-time feasibility can become difficult for fast systems or large models.
\end{itemize}

\subsection{MPC in CCD}

In CCD, the plant can be optimized together with MPC tuning parameters such as prediction horizon length, control horizon length, weighting matrices, or estimator settings. Such a formulation can be more realistic than pure OLOC because it uses a controller that can actually operate in closed loop.

\section{Moving from OLOC Solutions to Implementable Controllers}

\subsection{Why the transition is needed}

A common workflow in research is:
\begin{enumerate}[leftmargin=1.5em]
  \item solve an OLOC problem first,
  \item learn from the resulting ideal control pattern,
  \item build a simpler implementable controller that captures the most important features,
  \item co-design the plant with that implementable controller.
\end{enumerate}

This is often a very productive strategy. Figure~\ref{fig:ch6-implementable} summarizes the idea.

\begin{figure}[H]
\centering
\includegraphics[width=0.95\textwidth]{figures/fig_chp6_8.png}
\caption{A practical workflow from ideal open-loop optimization to an implementable closed-loop controller.}
\label{fig:ch6-implementable}
\end{figure}

\subsection{What can be learned from OLOC}

An OLOC solution may reveal:
\begin{itemize}[leftmargin=1.4em]
  \item when the control action tends to be active,
  \item whether the control behaves like damping, stiffness, or phase lead,
  \item what bandwidth or authority is needed,
  \item what physical plant characteristics are most valuable.
\end{itemize}

These observations can guide controller-structure selection.

\subsection{Typical transition paths}

Some common transitions are:
\begin{itemize}[leftmargin=1.4em]
  \item OLOC $\rightarrow$ PID or gain-scheduled feedback controller,
  \item OLOC $\rightarrow$ linear quadratic regulator (LQR) or state feedback,
  \item OLOC $\rightarrow$ MPC,
  \item OLOC $\rightarrow$ rule-based controller informed by optimal behavior,
  \item OLOC $\rightarrow$ parameterized controller identified from optimal trajectories.
\end{itemize}

\subsection{Why not skip OLOC entirely?}

Because OLOC can still be extremely useful. It often provides:
\begin{itemize}[leftmargin=1.4em]
  \item a best-case reference,
  \item insight into what control should do ideally,
  \item a benchmark against which practical controllers can be judged,
  \item a way to understand whether poor performance is caused by the plant or by controller limitations.
\end{itemize}

Thus, OLOC is often best viewed as an \emph{intermediate design tool} rather than as the final controller.

\section{Worked Examples}

\subsection{Worked Example 6.1: active suspension}

Suppose an active suspension is subjected to a known short road disturbance in simulation.

\textbf{Open-loop formulation:}
optimize actuator force $u(t)$ over the disturbance interval to minimize body acceleration and suspension deflection.

\textbf{Interpretation:}
this provides a best-case force history for that specific road input.

\textbf{Closed-loop implementation:}
construct a feedback controller based on measured suspension deflection and body velocity, then tune or co-design its gains.

The open-loop solution may show, for example, that strong control action is mainly needed at the onset of the disturbance and again when the suspension rebounds.

\subsection{Worked Example 6.2: wind turbine}

A wind-turbine study may optimize a blade-pitch trajectory against a known wind profile. This can indicate the best achievable load reduction for that profile, but the real turbine does not know the exact future gust history. A practical controller therefore uses measured rotor speed and other signals, possibly with preview or short-term forecasting. MPC or gain-scheduled feedback becomes a more realistic implementation path.

\subsection{Worked Example 6.3: marine energy device}

For a wave-energy converter, OLOC may optimize PTO force for a recorded wave sequence, giving a theoretical energy-capture benchmark. A practical controller must instead operate on measured or forecast wave information, perhaps using feedback or predictive control. Closed-loop CCD then redesigns the PTO or geometry together with that realizable controller.

\section{Comparison Table}

Table~\ref{tab:ch6-compare} summarizes the most important distinctions developed in this chapter.

\begin{table}[H]
\centering
\caption{Comparison of control representations relevant to CCD.}
\label{tab:ch6-compare}
\begin{tabular}{>{\raggedright\arraybackslash}p{2.9cm} >{\raggedright\arraybackslash}p{3.0cm} >{\raggedright\arraybackslash}p{3.2cm} >{\raggedright\arraybackslash}p{3.0cm}}
\toprule
Approach & Main decision variable & Main strength & Main limitation \\
\midrule
Open-loop optimal control & Full control trajectory $\uvec(t)$ or sampled control values & High flexibility and strong benchmark value & May require unrealistic future information and may not be directly implementable \\
Controller-parameter optimization & Gain or coefficient vector $\xc$ & Directly implementable and naturally feedback-based & Performance depends on chosen controller structure \\
Feedback-controller co-design & Plant variables plus controller parameters & Realistic closed-loop design & Can be more restrictive than full trajectory optimization \\
Model predictive control & Repeated finite-horizon optimization with updated state estimates & Closed-loop behavior with constraint handling and prediction & Requires online optimization and a good model \\
\bottomrule
\end{tabular}
\end{table}

\section{Common Mistakes and Misunderstandings}

\begin{enumerate}[leftmargin=1.5em]
  \item \textbf{Confusing a control trajectory with a controller.} A planned signal $u(t)$ is not the same thing as a feedback law.
  \item \textbf{Treating OLOC as automatically implementable.} An optimal open-loop signal may rely on impossible future information.
  \item \textbf{Ignoring information availability.} A controller can only use the information it can realistically measure, estimate, or forecast.
  \item \textbf{Assuming feedback always removes all uncertainty effects.} Feedback helps, but it does not eliminate model mismatch or measurement noise.
  \item \textbf{Thinking MPC is purely open loop.} Each MPC subproblem is open loop, but the overall strategy is closed loop because it is re-solved using updated measurements.
  \item \textbf{Assuming controller-parameter optimization is always inferior.} In practice, a well-chosen feedback structure can outperform a theoretically ideal but nonimplementable open-loop strategy on real uncertain systems.
\end{enumerate}

\section{Chapter Summary}

This chapter developed the distinction between open-loop optimal control and implementable closed-loop control in the context of CCD. We saw that open-loop optimal control directly optimizes a control trajectory and is valuable for performance benchmarks, insight, and idealized analysis. However, it often assumes future information that a real system does not possess.

We then distinguished trajectory optimization from controller-parameter optimization. Practical closed-loop systems generally use a controller law that maps available measurements into control actions. This viewpoint leads naturally to feedback-controller co-design, where plant variables and controller parameters are optimized together. We also introduced model predictive control as an important bridge: it repeatedly solves a finite-horizon open-loop problem, yet operates in closed loop by updating the optimization as new information becomes available.

Finally, we emphasized how to move from ideal open-loop solutions to implementable controllers. In practice, OLOC is often best used as a benchmark or design guide, while the final engineering design is based on a realizable feedback controller.

\section{Key Terms}

\begin{multicols}{2}
\begin{itemize}[leftmargin=1.4em,itemsep=0.25em]
  \item open-loop optimal control (OLOC)
  \item control trajectory
  \item controller parameterization
  \item feedback controller
  \item feedback-controller co-design
  \item future information
  \item information availability
  \item measured state
  \item estimated state
  \item forecast
  \item implementable controller
  \item model predictive control (MPC)
  \item receding horizon
  \item prediction horizon
  \item benchmark solution
  \item realizability
\end{itemize}
\end{multicols}

\section{Problems}

\subsection*{Conceptual problems}
\begin{enumerate}[leftmargin=1.5em,label=\textbf{P\arabic*.}]
  \item Define open-loop optimal control in your own words.
  \item What is the difference between a control trajectory and a controller?
  \item Why is feedback usually required in practical engineering systems?
  \item What does ``perfect future information'' mean in an OLOC setting?
  \item Why is MPC considered a feedback method even though it solves an open-loop problem internally?
\end{enumerate}

\subsection*{Comparison problems}
\begin{enumerate}[leftmargin=1.5em,label=\textbf{P\arabic*.},resume]
  \item Compare open-loop control and feedback control in terms of information usage.
  \item Compare trajectory optimization and controller-parameter optimization.
  \item Compare an ideal OLOC solution and an MPC implementation.
  \item Give one advantage and one limitation of OLOC in CCD.
  \item Give one advantage and one limitation of feedback-controller co-design.
\end{enumerate}

\subsection*{Formulation problems}
\begin{enumerate}[leftmargin=1.5em,label=\textbf{P\arabic*.},resume]
  \item Write a generic open-loop optimal-control problem with state dynamics and path constraints.
  \item Write a generic closed-loop control relation $\uvec=\kappa(\yvec;\xc)$ and explain each symbol.
  \item Write a generic feedback-controller co-design optimization problem involving plant variables $\xp$ and controller parameters $\xc$.
  \item In an MPC setting, list the sequence of steps performed at one time instant.
  \item Describe how the decision variables change when moving from OLOC to feedback-controller co-design.
\end{enumerate}

\subsection*{Interpretation problems}
\begin{enumerate}[leftmargin=1.5em,label=\textbf{P\arabic*.},resume]
  \item A student says, ``If the open-loop optimum is best, then that is the controller we should implement.'' Explain what is wrong with this statement.
  \item Another student says, ``MPC is not feedback because it computes an open-loop plan.'' Explain why this statement is incomplete.
  \item Why can the best plant under open-loop control differ from the best plant under feedback control?
  \item How can forecast uncertainty reduce the usefulness of an open-loop solution?
  \item Why can plant redesign sometimes improve information availability to the controller?
\end{enumerate}

\subsection*{Application problems}
\begin{enumerate}[leftmargin=1.5em,label=\textbf{P\arabic*.},resume]
  \item For an active suspension, explain how you would use OLOC as a preliminary design study before designing a feedback controller.
  \item For a wind-turbine problem, give one example of information that is measured and one example of information that is only forecast.
  \item For a marine-energy system, explain why a wave-record-based OLOC problem may be useful but not sufficient for final controller implementation.
  \item Choose an engineering system of your interest and explain whether you would prefer a direct feedback-controller co-design or an OLOC-first workflow.
  \item Give one example in which perfect future information is almost available and one example in which it is clearly unavailable.
\end{enumerate}

\subsection*{Mini-project problems}
\begin{enumerate}[leftmargin=1.5em,label=\textbf{P\arabic*.},resume]
  \item Pick an engineering system and draw two block diagrams: one for an open-loop optimal-control design study and one for an implementable feedback-controller design.
  \item Build a small comparison table for your system showing what information is available now, what can be predicted, and what remains unknown.
  \item Propose an MPC formulation for your chosen system, including the states, inputs, constraints, and prediction horizon.
  \item Describe how you would convert an ideal open-loop control history into a simpler practical controller structure.
  \item Write a one-page memo explaining whether OLOC should be viewed as a final controller or as a design-analysis tool in your chosen application.
\end{enumerate}

\section{Figure Files Supplied with This Chapter}

The following editable figure files are supplied separately as SVG files:
\begin{itemize}[leftmargin=1.4em]
  \item \texttt{fig\_chp6\_1.svg}
  \item \texttt{fig\_chp6\_2.svg}
  \item \texttt{fig\_chp6\_3.svg}
  \item \texttt{fig\_chp6\_4.svg}
  \item \texttt{fig\_chp6\_5.svg}
  \item \texttt{fig\_chp6\_6.svg}
  \item \texttt{fig\_chp6\_7.svg}
  \item \texttt{fig\_chp6\_8.svg}
\end{itemize}

\section{References and Further Reading}

\begin{thebibliography}{99}

\bibitem{betts}
J. T. Betts, \emph{Practical Methods for Optimal Control and Estimation Using Nonlinear Programming}. 2nd ed., SIAM, 2010.

\bibitem{brysonho}
A. E. Bryson and Y.-C. Ho, \emph{Applied Optimal Control}. Hemisphere, 1975.

\bibitem{allison2014aiaa}
J. T. Allison and D. R. Herber, ``Multidisciplinary design optimization of dynamic engineering systems,'' \emph{AIAA Journal}, vol. 52, no. 4, pp. 691--710, 2014.

\bibitem{allison2014jmd}
J. T. Allison, T. Guo, and Z. Han, ``Co-design of an active suspension using simultaneous dynamic optimization,'' \emph{Journal of Mechanical Design}, vol. 136, no. 8, 081003, 2014.

\bibitem{herber2019}
D. R. Herber and J. T. Allison, ``Nested and simultaneous solution strategies for general combined plant and control design problems,'' \emph{Journal of Mechanical Design}, vol. 141, no. 1, 011402, 2019.

\bibitem{garciasanz2019}
M. Garcia-Sanz, ``Control co-design: An engineering game changer,'' \emph{Advanced Control for Applications}, vol. 1, no. 1, 2019.

\bibitem{bayatbridging}
S. Bayat, L. Yang, and L. Zuo, ``Bridging the gap between open-loop and closed-loop control in co-design: A framework for complete optimal plant and control architecture design,'' 2026.

\bibitem{mayneMPC}
D. Q. Mayne, J. B. Rawlings, C. V. Rao, and P. O. M. Scokaert, ``Constrained model predictive control: Stability and optimality,'' \emph{Automatica}, vol. 36, no. 6, pp. 789--814, 2000.

\end{thebibliography}

\vfill
\begin{center}
\small\color{CCDgray}
End of Chapter 6
\end{center}

\end{document}
